\subsection{并查集（DSU）与 DSU on next 技巧}

朴素并查集（\texttt{find} 用\textbf{路径减半}，一行、非递归），外加 \texttt{taglr(l, r)} 把 $[l, r]$ 这段区间整段标记为「已处理」——具体做法是让区间里的每个点都不指向自己（\texttt{fa[i] = i + 1}，一路连到 $r+1$）。\texttt{find(x)} 因此天然返回 \texttt{x} 之后第一个未被标记的位置，扫描时常用来\textbf{跳过整段已处理区间}。\texttt{fa} 多开两格，\texttt{n + base} 是越界哨兵（全标完时 \texttt{find} 落在这里）。普通 \texttt{merge} 与 \texttt{taglr} 不要混在同一个对象上用（$l$、$l+1$ 已同块时会成环）。

\begin{minted}{cpp}
/**   并查集（DSU）+ DSU on next：find(x) = x 之后第一个没被标记的位置
 *    base = 0 / 1：合法下标 [0, n) / [1, n]；n 要给准，否则 countFa 出错
 */
struct DSU {
    vector<int> fa; int n = 0, base;
    DSU(int n = 0, int base = 0) : base(base) { init(n); }
    void init(int _n) { n = _n; fa.resize(n + base + 2); iota(all(fa), 0); }
    int find(int x) { while (x != fa[x]) x = fa[x] = fa[fa[x]]; return x; } // 路径减半
    bool same(int x, int y) { return find(x) == find(y); }
    bool merge(int x, int y) {          // 方向 fa[x] = y，配合 dsu on next
        x = find(x), y = find(y);
        if (x == y) return false;
        fa[x] = y; return true;
    }
    void taglr(int l, int r) {          // [l, r] 整段标记，之后 find 会越过去
        r = min(r, (int)SZ(fa) - 2);
        for (l = find(max(l, base)); l <= r; l = find(l)) fa[l] = l + 1;
    }
    int countFa() { int c = 0; for (int i = base; i < base + n; i++) c += fa[i] == i; return c; }
};
// 对拍：与旧版（ll + 两趟压缩）随机 2 万组一致
\end{minted}
